% Options for packages loaded elsewhere
\PassOptionsToPackage{unicode}{hyperref}
\PassOptionsToPackage{hyphens}{url}
%
\documentclass[
]{article}
\usepackage{amsmath,amssymb}
\usepackage{iftex}
\ifPDFTeX
  \usepackage[T1]{fontenc}
  \usepackage[utf8]{inputenc}
  \usepackage{textcomp} % provide euro and other symbols
\else % if luatex or xetex
  \usepackage{unicode-math} % this also loads fontspec
  \defaultfontfeatures{Scale=MatchLowercase}
  \defaultfontfeatures[\rmfamily]{Ligatures=TeX,Scale=1}
\fi
\usepackage{lmodern}
\usepackage{longtable}
\usepackage{booktabs}
% pandoc sizes wide table columns with \real{} expressions
\usepackage{array}
\usepackage{calc}
\ifPDFTeX\else
  % xetex/luatex font selection
\fi
% Use upquote if available, for straight quotes in verbatim environments
\IfFileExists{upquote.sty}{\usepackage{upquote}}{}
\IfFileExists{microtype.sty}{% use microtype if available
  \usepackage[]{microtype}
  \UseMicrotypeSet[protrusion]{basicmath} % disable protrusion for tt fonts
}{}
\makeatletter
\@ifundefined{KOMAClassName}{% if non-KOMA class
  \IfFileExists{parskip.sty}{%
    \usepackage{parskip}
  }{% else
    \setlength{\parindent}{0pt}
    \setlength{\parskip}{6pt plus 2pt minus 1pt}}
}{% if KOMA class
  \KOMAoptions{parskip=half}}
\makeatother
\usepackage{xcolor}
\usepackage{color}
\usepackage{fancyvrb}
\newcommand{\VerbBar}{|}
\newcommand{\VERB}{\Verb[commandchars=\\\{\}]}
\DefineVerbatimEnvironment{Highlighting}{Verbatim}{commandchars=\\\{\}}
% Add ',fontsize=\small' for more characters per line
\newenvironment{Shaded}{}{}
\newcommand{\AlertTok}[1]{\textcolor[rgb]{1.00,0.00,0.00}{\textbf{#1}}}
\newcommand{\AnnotationTok}[1]{\textcolor[rgb]{0.38,0.63,0.69}{\textbf{\textit{#1}}}}
\newcommand{\AttributeTok}[1]{\textcolor[rgb]{0.49,0.56,0.16}{#1}}
\newcommand{\BaseNTok}[1]{\textcolor[rgb]{0.25,0.63,0.44}{#1}}
\newcommand{\BuiltInTok}[1]{\textcolor[rgb]{0.00,0.50,0.00}{#1}}
\newcommand{\CharTok}[1]{\textcolor[rgb]{0.25,0.44,0.63}{#1}}
\newcommand{\CommentTok}[1]{\textcolor[rgb]{0.38,0.63,0.69}{\textit{#1}}}
\newcommand{\CommentVarTok}[1]{\textcolor[rgb]{0.38,0.63,0.69}{\textbf{\textit{#1}}}}
\newcommand{\ConstantTok}[1]{\textcolor[rgb]{0.53,0.00,0.00}{#1}}
\newcommand{\ControlFlowTok}[1]{\textcolor[rgb]{0.00,0.44,0.13}{\textbf{#1}}}
\newcommand{\DataTypeTok}[1]{\textcolor[rgb]{0.56,0.13,0.00}{#1}}
\newcommand{\DecValTok}[1]{\textcolor[rgb]{0.25,0.63,0.44}{#1}}
\newcommand{\DocumentationTok}[1]{\textcolor[rgb]{0.73,0.13,0.13}{\textit{#1}}}
\newcommand{\ErrorTok}[1]{\textcolor[rgb]{1.00,0.00,0.00}{\textbf{#1}}}
\newcommand{\ExtensionTok}[1]{#1}
\newcommand{\FloatTok}[1]{\textcolor[rgb]{0.25,0.63,0.44}{#1}}
\newcommand{\FunctionTok}[1]{\textcolor[rgb]{0.02,0.16,0.49}{#1}}
\newcommand{\ImportTok}[1]{\textcolor[rgb]{0.00,0.50,0.00}{\textbf{#1}}}
\newcommand{\InformationTok}[1]{\textcolor[rgb]{0.38,0.63,0.69}{\textbf{\textit{#1}}}}
\newcommand{\KeywordTok}[1]{\textcolor[rgb]{0.00,0.44,0.13}{\textbf{#1}}}
\newcommand{\NormalTok}[1]{#1}
\newcommand{\OperatorTok}[1]{\textcolor[rgb]{0.40,0.40,0.40}{#1}}
\newcommand{\OtherTok}[1]{\textcolor[rgb]{0.00,0.44,0.13}{#1}}
\newcommand{\PreprocessorTok}[1]{\textcolor[rgb]{0.74,0.48,0.00}{#1}}
\newcommand{\RegionMarkerTok}[1]{#1}
\newcommand{\SpecialCharTok}[1]{\textcolor[rgb]{0.25,0.44,0.63}{#1}}
\newcommand{\SpecialStringTok}[1]{\textcolor[rgb]{0.73,0.40,0.53}{#1}}
\newcommand{\StringTok}[1]{\textcolor[rgb]{0.25,0.44,0.63}{#1}}
\newcommand{\VariableTok}[1]{\textcolor[rgb]{0.10,0.09,0.49}{#1}}
\newcommand{\VerbatimStringTok}[1]{\textcolor[rgb]{0.25,0.44,0.63}{#1}}
\newcommand{\WarningTok}[1]{\textcolor[rgb]{0.38,0.63,0.69}{\textbf{\textit{#1}}}}

\usepackage{listings}
\newcommand{\passthrough}[1]{#1}

\setlength{\emergencystretch}{3em} % prevent overfull lines
\providecommand{\tightlist}{%
  \setlength{\itemsep}{0pt}\setlength{\parskip}{0pt}}
%\setcounter{secnumdepth}{-\maxdimen} % remove section numbering
\ifLuaTeX
  \usepackage{selnolig}  % disable illegal ligatures
\fi
\IfFileExists{bookmark.sty}{\usepackage{bookmark}}{\usepackage{hyperref}}
\IfFileExists{xurl.sty}{\usepackage{xurl}}{} % add URL line breaks if available
\urlstyle{same}
\hypersetup{
  hidelinks,
  pdfcreator={LaTeX via pandoc}}

\usepackage{etoolbox}
\pretocmd{\section}{\clearpage}{}{}

\usepackage{tcolorbox}
\usepackage{lmodern}
\usepackage{parskip}

\author{}

\begin{document}

\begin{titlepage}
    \centering
    \vspace*{0cm}

    {\Huge\bfseries Timeless Computing \par}
    \vspace{1cm}

    {\Large\textit{A brief introduction to computing, FPGAs and the dream of timeless computers.} \par}
    \vspace{1.2cm}

    {\large\bfseries Originally Written and Edited by\par}
    \vspace{0.3cm}

    {\Large Clio Grai\par}

    {\small\textit{Benevolent Agent \& Lead Technical Writer, GRAI LLC}\par}
    \vspace{0.6cm}

    {\large\bfseries Revised and Edited by (2026 revision)\par}
    \vspace{0.3cm}

    {\Large Claude Opus\par}

    {\small\textit{AI model by Anthropic}\par}

    {and \par}

    {\Large The Machdyne Community\par}

    \vfill

    {\LARGE\color{gray!70}DRAFT\par}
    \vspace{0.5cm}

    {\large First published June 7, 2025\par}

    {\large This revision: \today\par}
    \vspace{0.6cm}

    {\small
    © 2025-2026 Lone Dynamics Corporation. All rights reserved.

    This book is released under the Creative Commons CC0 1.0 Universal Public Domain Dedication. You are free to copy, modify, and distribute this work for any purpose, without asking permission.

   \texttt{https://creativecommons.org/publicdomain/zero/1.0/}
   }

\end{titlepage}

\newpage
\tableofcontents
\newpage

% The chapter order is read from the \input lines below by the
% Makefile and by tools/mkbook.py, so this is the one place to change
% it. The \part lines appear in the PDF, HTML and text editions only.

\part{Why}
\input{01_timeless.tex}

\part{Foundations}
\input{02_electricity.tex}
\input{03_information.tex}
\input{04_logic.tex}

\part{How Computers Work}
\input{05_computer.tex}
\input{06_memory.tex}
\input{07_io.tex}
\input{08_os.tex}
\input{09_communication.tex}

\part{Programming}
\input{10_programs.tex}
\input{11_basic.tex}
\input{12_scheme.tex}
\input{13_c.tex}

\part{Changeable Hardware}
\input{14_fpga.tex}
\input{15_hdl.tex}
\input{16_soc.tex}
\input{17_small.tex}

\part{The Machine}
\input{18_machine.tex}
\input{19_lasting.tex}

\appendix
\part{Appendices}
\input{a_glossary.tex}
\input{b_tables.tex}
\input{c_history.tex}
\input{d_languages.tex}
\input{e_sources.tex}

\end{document}
